\clearpage\SourceLeaf{366}\ArithmeticHead{359}{R}{ZUR THEORIE DER REIHE $F(\alpha,\beta,\gamma,x)$.\quad IV.}
\NumberedItem{6}
\noindent\hspace*{\parindent}Nachdem dieser Satz fest steht, beachte man zunächst, dass die Entwickelung von
\[
(1-ff)^{-\theta-\mu}
\]
\noindent lauter Glieder von einerlei Zeichen hat, sobald $\mu+\theta$ positiv geworden ist. Dann ergibt sich aus dem Hülfssatz leicht, dass
\[
F(\theta,n+\theta,n+1,ff)=A
\]
\noindent gesetzt, [ferner]
\[
(1-ff)^{-\theta}F\left(\theta,1-\theta,n+1,-\frac{ff}{1-ff}\right)=B_\mu,
\]
\noindent insofern von $F\left(\theta,1+\theta,n+1,-\frac{ff}{1-ff}\right)$ die ersten $\mu$ Glieder summirt werden, während das nächstfolgende Glied $b_{\mu+1}$ heisst, $A-B_\mu$ dasselbe Zeichen haben wird, wie $b_{\mu+1}$.

\indent Jener Hülfssatz ist übrigens ein spezieller Fall von dem Allgemeinern:

\indent Wenn $a,b,c,d$ u.s.w. positive Grössen sind, jede folgende grösser als die vorhergehende, so ist die Summe
\[
a-kb+\frac{k\ k-1}{1{.}\,2}c-\frac{k{.}\,k-1{.}\,k-2}{1{.}\,2{.}\,3}d+\text{u.s.w.}
\]
\noindent wo man auch abbricht, immer mit demselben Zeichen behaftet, wie das letzte zugezogene Glied.

\indent Man wendet diess an, indem man von hinten anfängt. Der Satz bleibt auch wahr, wenn von irgend einer Stelle an die Grössen $a,b,c,d$ u.s.w. alternirende Zeichen haben.
\par\bigskip\begin{center}\bfseries BEMERKUNG.\end{center}
{\small
\noindent\hspace*{\parindent}Dieser Aufzeichnung (der letzten des Handbuchs 19, Be) geht unmittelbar vorher (auf S. 288 des Handbuchs) die von \textsc{Schering} in den \textls[120]{Bemerkungen} Werke III, S. 230 erwähnte Berechnung von $\log\Pi i$ und $\Pi i$. Daß diese, wie \textsc{Schering} angibt, »\textls[120]{nach dem Jahre 1847}« anzusetzen ist, folgt daraus, daß sich auf S. 276 des Handbuchs die Zeitangabe »1847, Mai 13« findet. Natürlich gilt dieselbe Zeitbestimmung für die vorstehend abgedruckte Aufzeichnung. Zu dem Gegenstande vergleiche man den art. 6. der \textit{Disquisitiones circa seriem}, Werke III, S. 128.
\par\smallskip\hfill\textsc{Schlesinger.}\par
}
\ClosingRule
\clearpage
